% META: {"id":"1SPE-DERIVATION-LOCAL-CR-013","chapitre":"1SPE-DERIVATION-LOCAL","type_objet":"cours","capacites_codes":["C4"],"sources_inspiration":[],"status":"generated"}
\section{Équation de la tangente}
\theoreme[Équation de la tangente]{Si $f$ admet un nombre dérivé en $a$, la tangente à la courbe de $f$ au point d'abscisse $a$ a pour équation
\[
y=f(a)+f'(a)(x-a).
\]
L'hypothèse de dérivabilité en $a$ est nécessaire : elle garantit que la pente $f'(a)$ existe.}
\exemple{Si $f(1)=3$ et $f'(1)=2$, la tangente a pour équation $y=3+2(x-1)$, soit $y=2x+1$.}
\contreexemple{L'équation $y=f'(a)x$ oublie en général l'ordonnée $f(a)$ : avec les données précédentes, $y=2x$ ne passe pas par $(1;3)$.}
\propriete[Vérification]{Dans l'équation de la tangente, remplacer $x$ par $a$ donne $y=f(a)$. La droite passe donc par le point de contact.}
\exemple{Pour $y=3+2(x-1)$, remplacer $x$ par $1$ donne $y=3$.}
\contreexemple{Une droite de pente $2$ passant par $(0;3)$, d'équation $y=2x+3$, n'est pas la tangente précédente : elle ne passe pas par $(1;3)$.}
\begin{nxfigure}{Tangente et point de contact}
\begin{tikzpicture}
\begin{axis}[nexus,
  xmin=-1, xmax=4.5, ymin=-1, ymax=6,
  xlabel=$x$, ylabel=$y$,
  xtick={0,1,2,3,4}, ytick={0,1,...,5}]
  % Courbe f(x) = 0.4x^3 - 1.5x^2 + 2x + 1
  \addplot[courbe,domain=-0.5:4,samples=80] {0.4*x^3 - 1.5*x^2 + 2*x + 1};
  % Tangente en a=1 : f(1) = 1.9, f'(1) = 1.2*1 - 3*1 + 2 = 0.2
  \addplot[tangente,domain=-0.5:3.5] {0.2*(x-1) + 1.9};
  % Point de contact
  \addplot[point] coordinates {(1,1.9)};
  
\node[anchor=south west,font=\footnotesize,nxBleu] at (axis cs:1,1.9) {$A(a;f(a))$};
  % Annotation tangente
  
\node[anchor=south,font=\footnotesize,nxOr] at (axis cs:3,2.3) {$y=f(a)+f'(a)(x-a)$};
\end{axis}
\end{tikzpicture}

\nxcaption{La tangente passe par $A$ avec la pente $f'(a)$.}
\end{nxfigure}

\margeAppui{Retenir les trois informations séparées : l'abscisse $a$, l'ordonnée $f(a)$ et la pente $f'(a)$.}
\margeAppui{La forme $y-f(a)=f'(a)(x-a)$ est la forme « point-pente » d'une droite ; on peut ensuite développer si nécessaire.}
\erreurFrequente{\textbf{Copie fautive :} $y=f(a)+f'(a)(x+a)$. \textbf{Correction :} l'écart horizontal se mesure par $x-a$, donc le signe est négatif.}
\erreurFrequente{\textbf{Copie fautive :} $y=f(a)+f'(x)(x-a)$. \textbf{Correction :} on cherche une droite : sa pente est le nombre fixe $f'(a)$, non une quantité dépendant de $x$.}
\approfondissement{\textbf{Démonstration exigible.} Toute droite de coefficient directeur $m$ passant par $A(a;f(a))$ vérifie $y-f(a)=m(x-a)$. En effet, pour tout point $M(x;y)$ de la droite distinct de $A$, son coefficient directeur est $\frac{y-f(a)}{x-a}=m$, d'où l'égalité ; elle reste vraie pour $x=a$. La tangente possède le coefficient directeur $m=f'(a)$. On obtient donc $y=f(a)+f'(a)(x-a)$. Réciproquement, la droite ainsi définie passe par $A$ et possède la pente $f'(a)$ : elle est bien la tangente.}

% BEGIN-VERIFY
% from sympy import symbols, diff, simplify, Rational
%
% x = symbols('x', real=True)
%
% def tangente(f, a):
%     """y = f(a) + f'(a)(x - a), construite depuis la fonction elle-meme."""
%     return f.subs(x, a) + diff(f, x).subs(x, a) * (x - a)
%
% # L'exemple redige : f(1) = 3 et f'(1) = 2 donnent y = 2x + 1.
% f = x**2 + 2*x            # f(1) = 3, f'(1) = 4 : on prend plutot une fonction
%                           # qui realise exactement les donnees de l'enonce.
% g = 2*x + 1               # f(1)=3, f'(1)=2
% assert g.subs(x, 1) == 3 and diff(g, x).subs(x, 1) == 2
% assert simplify(tangente(g, 1) - (2*x + 1)) == 0
%
% # La propriete de verification : en x = a, l'equation redonne f(a).
% for fonction, a in ((x**2, 3), (x**3 - x, -1), (2*x + 1, 1)):
%     assert simplify(tangente(fonction, a).subs(x, a) - fonction.subs(x, a)) == 0
%
% # Le contre-exemple : y = f'(a)x oublie l'ordonnee. Avec f(1)=3 et f'(1)=2,
% # la droite y = 2x ne passe pas par (1;3).
% assert (2 * 1) != 3
% # Et la droite y = 2x + 3, de meme pente, n'y passe pas non plus.
% assert (2 * 1 + 3) != 3
%
% # La tangente de la figure : f(x) = 0.4x^3 - 1.5x^2 + 2x + 1 en a = 1.
% courbe = Rational(2, 5)*x**3 - Rational(3, 2)*x**2 + 2*x + 1
% assert courbe.subs(x, 1) == Rational(19, 10)
% assert diff(courbe, x).subs(x, 1) == Rational(1, 5)
% assert simplify(tangente(courbe, 1)
%                 - (Rational(1, 5)*(x - 1) + Rational(19, 10))) == 0
% END-VERIFY
